\specsection{Exercises for the Annex}

\begin{enumerate}
\def\labelenumi{\arabic{enumi})}
\tightlist
\item
  Let E be a finite-dimensional vector space over a commutative field of characteristic \(\neq 2\) . Denote by H a nondegenerate quadratic form on E, and by S the symmetric bilinear form on \(E \times E\) such that \(\mathrm{H}(x) = \mathrm{S}(x, x)\) for all x in E. Let Q be a quadratic form on E.
\end{enumerate}

\begin{enumerate}
\def\labelenumi{\alph{enumi})}
\item
  There exists an endomorphism \(u\) of \(\mathbf{E}\) characterized by the relations \(\mathbf{Q}(x) = \mathrm{S}(u(x), x)\) and \(\mathrm{S}(u(x), y) = \mathrm{S}(x, u(y))\) for all \(x\) and \(y\) in \(\mathbf{E}\) . One sets \(\operatorname{Tr}(\mathbf{Q}/\mathbf{H}) = \operatorname{Tr}(u)\) .
\item
  Generalize Remark 3 of No.~1.
\item
  If \((e_i)_{1\leqslant i\leqslant m}\) is a basis of E orthonormal for H, then \(\operatorname{Tr}(\mathbf{Q}/\mathbf{H}) = \sum_{i=1}^{m}\mathrm{Q}(e_i)\) .
\end{enumerate}

\begin{enumerate}
\def\labelenumi{\arabic{enumi})}
\setcounter{enumi}{1}
\tightlist
\item
  Let E be a real Hilbert space, Q a continuous positive quadratic form on E, and H the quadratic form \(x \mapsto \|x\|^2\) on E.
\end{enumerate}

\begin{enumerate}
\def\labelenumi{\alph{enumi})}
\item
  Show that \(\operatorname{Tr}(\mathbf{Q}/\mathbf{H}) = \sum_{i \in I} \mathbf{Q}(e_i)\) for every orthonormal basis \((e_i)_{i \in I}\) of E. (Let \((a_1, \ldots, a_p)\) be a finite orthonormal family in E; for every \(\varepsilon > 0\) , there exist a finite subset J of I and elements \(a_1', \ldots, a_p'\) that are linear combinations of the \(e_i\) for \(i \in J\) and are such that \(\|a_j - a_j' \| < \varepsilon\) for \(1 \leqslant j \leqslant p\) ; then \(\sum_{j=1}^{p} \mathbf{Q}(a_j') \leqslant \sum_{i \in J} \mathbf{Q}(e_j) \leqslant \sum_{i \in I} \mathbf{Q}(e_i)\) . Deduce from this that \(\sum_{j=1}^{p} \mathbf{Q}(a_j) \leqslant \sum_{i \in I} \mathbf{Q}(e_i)\) .)
\item
  Deduce from \(a\) ) a new proof of Prop. 3.
\item
  Let \(\mathbf{E}_0\) be a dense linear subspace of \(\mathbf{E}\) , \(\mathbf{Q}_0\) (resp. \(\mathbf{H}_0\) ) the restriction of \(\mathbf{Q}\) (resp. \(\mathbf{H}\) ) to \(\mathbf{E}_0\) . Show that \(\operatorname{Tr}(\mathbf{Q}/\mathbf{H}) = \operatorname{Tr}(\mathbf{Q}_0/\mathbf{H}_0)\) . (Let \((e_1, \ldots, e_n)\) be a linearly independent sequence in \(\mathbf{E}\) , generating a subspace \(\mathbf{F}\) . Denote by \(\mathbf{Q}_{\mathbf{F}}\) (resp. \(\mathbf{H}_{\mathbf{F}}\) ) the restriction of \(\mathbf{Q}\) (resp. \(\mathbf{H}\) ) to \(\mathbf{F}\) . Using Remark 3 of No.~1, show that \(\operatorname{Tr}(\mathbf{Q}_{\mathbf{F}}/\mathbf{H}_{\mathbf{F}})\) is a continuous function of \((e_1, \ldots, e_n)\) .)
\end{enumerate}

\begin{enumerate}
\def\labelenumi{\arabic{enumi})}
\setcounter{enumi}{2}
\item
  Let E and F be two real vector spaces and u a linear mapping of E into F; if Q and H are positive quadratic forms on F such that \(\mathrm{H}(x)=0\) implies \(\mathrm{Q}(x)=0\) for \(x\in F\) , then \(\mathrm{Tr}\left(\mathrm{Q}\circ\mathrm{u}/\mathrm{H}\circ\mathrm{u}\right)\leqslant\mathrm{Tr}\left(\mathrm{Q}/\mathrm{H}\right)\) .
\item
  Let I be a set and \(l^2(\mathbf{I})\) the vector space of families \(\mathbf{x} = (x_i)_{i \in \mathbf{I}}\) of real numbers such that \(\sum_{i \in \mathbf{I}} x_i^2\) is finite. Let \((\lambda_i)_{i \in \mathbf{I}}\) be a summable family of positive numbers. For every \(\mathbf{x}\) in \(l^2(\mathbf{I})\) , set \(Q(\mathbf{x}) = \sum_{i \in \mathbf{I}} \lambda_i x_i^2\) and \(H(\mathbf{x}) = \sum_{i \in \mathbf{I}} x_i^2\) . Show that \(\operatorname{Tr}(Q / H) = \sum_{i \in \mathbf{I}} \lambda_i\) . (Make use of Exercise 2 a).
\item
  Let E be a real vector space, \((\lambda_i)_{i \in I}\) a summable family of positive numbers, and \((y_i)_{i \in I}\) a family of linear forms on E. For every \(x \in E\) , set \(\mathrm{H}(x) = \sum_{i \in I} \langle x, y_i \rangle^2\) and \(\mathrm{Q}(x) = \sum_{i \in I} \lambda_i \langle x, y_i \rangle^2\) ; assume that \(\mathrm{H}(x)\) is finite for every \(x \in E\) . Show that \(\mathrm{Q}(x)\) is finite for every \(x \in E\) , that Q and H are positive quadratic forms on E, and that \(\operatorname{Tr}(\mathrm{Q}/\mathrm{H}) \leqslant \sum_{i \in I} \lambda_i\) . (Set \(u(x) = (\langle x, y_i \rangle)_{i \in I}\) and apply Exer. 3 to the linear mapping \(u\) of E into \(l_2(I)\) .)
\item
  Let E be a real vector space, Q, H and \(H_{0}\) positive quadratic forms on E. Assume that \(H \leqslant a \cdot H_{0}\) , where a is a positive real number. Prove that \(\operatorname{Tr}(Q/H_{0}) \leqslant a \cdot \operatorname{Tr}(Q/H)\) . (Reduce to the case that E is finite-dimensional by Remark 1 of No.~1, then conclude by Prop. 1, using the existence of a basis of E that is orthogonal for both H and \(H_{0}\) .)
\item
  Let E be a real Hilbert space. Show that the Sazonov topology on E is defined by the set of semi-norms \(Q^{1/2}\) , where Q runs over the set of nuclear positive quadratic forms on E. (Use Exer. 6.)
\item
  Let E and F be two real Hilbert spaces. There exists on \(E \otimes F\) a Hausdorff pre-Hilbert space structure such that \((x_{1} \otimes y_{1}|x_{2} \otimes y_{2}) = (x_{1}|x_{2}) \cdot (y_{1}|y_{2})\) for \(x_{1}, x_{2}\) in E and \(y_{1}, y_{2}\) in F. One denotes by \(E \otimes_{2} F\) the Hilbert space completion of \(E \otimes F\) .
\end{enumerate}

\begin{enumerate}
\def\labelenumi{\alph{enumi})}
\item
  If \(\mathbf{E}'\) (resp. \(\mathbf{F}'\) ) is a closed linear subspace of \(\mathbf{E}\) (resp. \(\mathbf{F}\) ), show that \(\mathbf{E}' \otimes_2 \mathbf{F}'\) may be identified with the closed linear subspace of \(\mathbf{E} \otimes_2 \mathbf{F}\) generated by the elements \(x \otimes y\) such that \(x \in \mathbf{E}'\) and \(y \in \mathbf{F}'\) .
\item
  Suppose that E (resp. F) is the hilbertian sum of a family \((\mathrm{E}_{\alpha})_{\alpha\in\mathbf{A}}\) (resp. \((\mathrm{F}_{\beta})_{\beta\in\mathbf{B}}\) ) of closed linear subspaces (TVS, V, §2, No.~2, Def. 2). Show that \(E\otimes_{2}F\) is the hilbertian sum of the family \((\mathrm{E}_{\alpha}\otimes_{2}\mathrm{F}_{\beta})_{(\alpha,\beta)\in\mathbf{A}\times\mathbf{B}}\) of closed linear subspaces.
\item
  If \((e_i)_{i\in \mathbf{I}}\) (resp. \((f_j)_{j\in \mathbf{J}})\) is an orthonormal basis of E (resp. F), then the family \((e_i\otimes f_j)_{(i,j)\in \mathbf{I}\times \mathbf{J}}\) is an orthonormal basis of \(\mathbf{E}\otimes_2\mathbf{F}\) .
\item
  Let \(\mathbf{G}\) be a real Hilbert space. Define canonical isomorphisms of \(\mathbf{E} \otimes_{2} \mathbf{F}\) onto \(\mathbf{F} \otimes_{2} \mathbf{E}\) and of \((\mathbf{E} \otimes_{2} \mathbf{F}) \otimes_{2} \mathbf{G}\) onto \(\mathbf{E} \otimes_{2} (\mathbf{F} \otimes_{2} \mathbf{G})\) .
\item
  Let \(E_{1}\) and \(F_{1}\) be two real Hilbert spaces, and \(u:E\to E_{1}, v:F\to F_{1}\) continuous linear mappings. Show that \(u\otimes v\) may be extended by continuity to a continuous linear mapping \(u\otimes_{2}v\) of \(E\otimes_{2}F\) into \(E_{1}\otimes_{2}F_{1}\) , and that \(\|u\otimes_{2}v\|=\|u\|\cdot\|v\|\) .
\end{enumerate}

\begin{enumerate}
\def\labelenumi{\arabic{enumi})}
\setcounter{enumi}{8}
\tightlist
\item
  Let E and F be two real Hilbert spaces.
\end{enumerate}

\begin{enumerate}
\def\labelenumi{\alph{enumi})}
\item
  Show that there exists a continuous linear mapping \(\varphi\) of \(\mathbf{E} \otimes_2 \mathbf{F}\) into \(\mathcal{L}(\mathbf{E}; \mathbf{F})\) characterized by \(\varphi(x \otimes y)(x') = (x|x') \cdot y\) for \(x, x'\) in \(\mathbf{E}\) and \(y\) in \(\mathbf{F}\) . Show that \(\varphi\) has norm 1, and that \(\varphi(\mathbf{E} \otimes \mathbf{F})\) is the set of continuous linear mappings of finite rank of \(\mathbf{E}\) into \(\mathbf{F}\) .
\item
  Show that \(\varphi\) is a bijection of \(E \otimes_2 F\) onto the set HS (E, F) of Hilbert--Schmidt linear mappings of E into F (make use of Exer. 8 c) and Prop. 3). Equip HS (E, F) with the Hilbert space structure deduced from that of \(E \otimes_2 F\) by means of the bijection \(\varphi\) ; the corresponding norm is denoted \(\|u\|_2\) . Let \(u \in \mathrm{HS}(E, F)\) ; one defines positive quadratic forms \(Q_u\) and H on E by \(Q_u(x) = \|u(x)\|^2\) and \(H(x) = \|x\|^2\) . Show that \(\|u\|_2^2 = \operatorname{Tr}(Q_u / H)\) .
\item
  Let \(\mathbf{E}_1\) and \(\mathbf{F}_1\) be two real Hilbert spaces, and \(u: \mathbf{E}_1 \to \mathbf{E}\) , \(v: \mathbf{F} \to \mathbf{F}_1\) continuous linear mappings. Let \(\varphi_1\) be the isomorphism of \(\mathbf{E}_1 \otimes_2 \mathbf{F}_1\) onto \(\mathrm{HS}(\mathbf{E}_1, \mathbf{F}_1)\) defined in a manner analogous to \(\varphi\) . Show that \(v \circ \varphi(t) \circ u = \varphi_1((u^* \otimes v)(t))\) for all \(t \in \mathbf{E} \otimes_2 \mathbf{F}\) . From this, deduce that for every Hilbert-Schmidt mapping \(w\) of \(\mathbf{E}\) into \(\mathbf{F}\) , the linear mapping \(v \circ w \circ u\) of \(\mathbf{E}_1\) into \(\mathbf{F}_1\) is Hilbert-Schmidt, and that \(\|v \circ w \circ u\|_2 \leqslant \|v\| \cdot \|w\|_2 \cdot \|u\|\) .
\end{enumerate}

